\documentclass[11pt]{article}

\usepackage{series}
\usepackage{booktabs}

\hypersetup{
  pdftitle={Dirac Delta Limits, Spectral Kernels, and the MTT Finite-Kernel Diagnostic},
  pdfauthor={Peter Nero},
  pdfsubject={Dirac distributions, spectral projectors, heat kernels, constraint tubes, gauge fixing, and Modal Triplet Theory},
  pdfkeywords={Dirac delta, spectral projector, heat kernel, approximate identity, coarea formula, Faddeev-Popov, Modal Triplet Theory},
  hidelinks
}

\newcommand{\Dcal}{\mathcal D}
\newcommand{\Hh}{\mathcal H}
\newcommand{\Jcal}{\mathcal J}
\newcommand{\Fock}{\mathcal F}
\newcommand{\Tr}{\operatorname{Tr}}
\newcommand{\diag}{\operatorname{diag}}

\title{Dirac Delta Limits, Spectral Kernels, and the MTT Finite-Kernel Diagnostic\\
\large Exact Approximate Identities, Constraint Tubes, and Scope Boundaries}

\author{Peter Nero}

\date{July 2026, Version 1}

\begin{document}
\maketitle

\begin{abstract}
Dirac delta distributions play several mathematically different roles in
physics. They are identity kernels, point sources, hard constraints,
conservation distributions, formal densities of continuous observables, and
local gauge-slice selectors. This paper separates those roles before asking
which admit finite-resolution representatives. On a compact elliptic
background, spectral projector kernels converge to the diagonal delta
distribution, with the Sobolev truncation estimate
\[
\|(I-\Pi_\Lambda)f\|_{H^s}
\leq (1+\Lambda)^{-r/2}\|f\|_{H^{s+r}}.
\]
Heat kernels give a different, nonprojective approximate identity and obey
\[
\|(e^{-\tau\Delta}-I)f\|_{H^s}
\leq \tau^{r/2}\|f\|_{H^{s+r}},
\qquad 0\leq r\leq2.
\]
For a smooth submersion \(C:\mathbb R^n\to\mathbb R^m\), normalized Gaussian
tubes converge by the coarea formula to the correctly Jacobian-weighted
constraint distribution on \(C^{-1}(0)\). Compression of the canonical
commutation relations by an orthogonal projection replaces the full inner
product kernel by the projected kernel exactly, but only in that declared
compressed theory. These results support a restricted Modal Triplet Theory
(MTT) diagnostic: when a delta occurs in a downstream model, one should ask
whether selected upstream geometry emits a finite projector, smoothing
kernel, source profile, detector response, or constraint tube, and in which
topology the sharp limit is controlled. They do not imply that every delta is
a hidden projection, that a spectral gap fixes a physical width, or that
finite smearing by itself solves gauge fixing, measurement, ultraviolet
renormalization, or outcome selection.
\end{abstract}

\section*{Version 1 Revision Note}
\addcontentsline{toc}{section}{Version 1 Revision Note}
\begin{description}
\item[Supersedes] The unversioned April 2026 manuscript
\emph{Dirac Delta Functions as Singular Shadows of Admissible Projection: A
Fixed-Point and Gauge-Theoretic Formulation in Modal Triplet Theory}.
\item[Reason] The earlier manuscript treated all physical delta distributions
as zero-width projections, inferred finite physical width from a spectral gap,
assigned a delta outcome to every deterministic fixed-point basin, and
extended a local Faddeev--Popov Jacobian argument to global functional gauge
fixing. It also used an unnormalized Gaussian constraint filter and suggested
that replacing deltas by finite kernels generally resolves renormalization.
\item[Resolution] This version classifies the distinct uses of the delta,
proves separate spectral, heat-kernel, constraint-tube, and compressed-CCR
statements with explicit hypotheses and error topologies, corrects the
fixed-point claim, and limits the gauge argument to a local finite-dimensional
model. Measurement and field-theory applications are stated at their actual
tiers.
\item[Retained result] Spectral projectors and heat kernels provide rigorous
finite kernels approaching the identity distribution, normalized constraint
tubes converge by coarea, and finite-sector compression yields an exact
projected commutator kernel.
\item[Remaining boundary] No theorem here derives a particular cutoff,
detector width, source profile, gauge slice, or regularization scale from the
selected MTT carrier. Global non-Abelian gauge fixing, arbitrary measurement
contexts, one-history completion, and symmetry-preserving ultraviolet
completion remain separate problems.
\end{description}

\tableofcontents

\section{Why the uses of \texorpdfstring{\(\delta\)}{delta} must be separated}

The Dirac delta is a distribution, not an ordinary function. On
\(\mathbb R^d\), its defining property is
\[
\langle\delta_{x_0},f\rangle=f(x_0),
\qquad f\in C_c^\infty(\mathbb R^d).
\]
That single definition supports several constructions whose physical meanings
are not interchangeable. A delta may represent an exact identity operator, an
ideal point source, a constraint surface with a Jacobian, or the Fourier
expression of an exact symmetry. In other settings it is only a formal density
for an operator-valued measure.

The earlier projection reading noticed something useful: many sharp
distributional objects can be approached by finite kernels. The mistake was
to turn that useful construction into a universal origin claim. The corrected
question is narrower:
\begin{quote}
For this particular occurrence of a delta, is there a selected finite object
whose controlled limit gives the distribution, and does replacing the delta
preserve the mathematical and physical structure that matters?
\end{quote}

\begin{center}
\begin{tabular}{p{0.23\linewidth}p{0.33\linewidth}p{0.34\linewidth}}
\toprule
Use & Exact object & Possible finite representative\\
\midrule
Identity kernel &
Diagonal delta distribution &
Spectral projector or heat kernel\\
Point source &
Distributional forcing term &
Extended source density\\
Hard constraint &
Level-set distribution with Jacobian &
Normalized constraint tube\\
Momentum conservation &
Fourier distribution from translation invariance &
Finite spacetime-window transform\\
PVM density &
Formal density of a PVM &
Detector-smeared POVM\\
Gauge slice &
Local constraint and orbit Jacobian &
Finite gauge tube, locally\\
\bottomrule
\end{tabular}
\end{center}

The third column is not automatic. Each proposed finite representative needs a
normalization, a convergence topology, and a source theorem if it is to be a
prediction rather than a regulator chosen by hand.

\section{The diagonal delta is the identity kernel}

Let \(X\) be a compact smooth Riemannian manifold with volume form
\(\dd V\). The diagonal delta distribution
\(\delta_{\mathrm{diag}}\in\Dcal'(X\times X)\) is defined by
\[
\langle\delta_{\mathrm{diag}},\Psi\rangle
=\int_X\Psi(x,x)\,\dd V(x),
\qquad \Psi\in C^\infty(X\times X).
\]
With the usual kernel convention, it represents the identity:
\[
f(x)=\int_X\delta_{\mathrm{diag}}(x,y)f(y)\,\dd V(y).
\]
This statement is exact. It does not assert that a finite physical process has
generated the identity operator.

An orthogonal projection \(P:\Hh\to\Hh\) is also an exact bounded operator.
If it has a kernel \(K_P\), that kernel represents the identity only on
\(\Ran(P)\):
\[
Pf=f\qquad\text{for }f\in\Ran(P).
\]
An arbitrary bounded projection need not possess a pointwise smooth or bounded
kernel. Kernel regularity follows only after additional hypotheses such as
finite rank, Hilbert--Schmidt regularity, or smoothing. This qualification is
essential whenever a finite coherent projector is presented as a physical
profile.

\section{Spectral projectors approach the identity}

\begin{assumption}[Compact elliptic setting]\label{ass:elliptic}
Let \(\Delta\geq0\) be a nonnegative self-adjoint Laplace-type operator on
\(L^2(X)\), where \(X\) is compact and has either no boundary or a fixed
self-adjoint elliptic boundary condition. Let
\[
\Delta\phi_j=\lambda_j\phi_j,
\qquad
0\leq\lambda_0\leq\lambda_1\leq\cdots,
\]
where \(\{\phi_j\}\) is a complete orthonormal eigenbasis.
\end{assumption}

For \(\Lambda\geq0\), define
\[
\Pi_\Lambda f
=\sum_{\lambda_j\leq\Lambda}
\langle\phi_j,f\rangle\phi_j
\]
and
\[
K_\Lambda(x,y)
=\sum_{\lambda_j\leq\Lambda}
\phi_j(x)\overline{\phi_j(y)}.
\]
Each \(\Pi_\Lambda\) is a finite-rank orthogonal projection and
\(K_\Lambda\) is smooth.

\begin{theorem}[Spectral projector kernel limit]
\label{thm:spectral-kernel}
Under Assumption~\ref{ass:elliptic},
\[
K_\Lambda\longrightarrow\delta_{\mathrm{diag}}
\quad\text{in }\Dcal'(X\times X)
\]
as \(\Lambda\to\infty\). Equivalently,
\(\Pi_\Lambda f\to f\) in \(C^\infty(X)\) for every
\(f\in C^\infty(X)\).
\end{theorem}

\begin{proof}
For \(s\in\mathbb R\), use the spectral Sobolev norm
\[
\|f\|_{H^s}^2
=\sum_j(1+\lambda_j)^s
|\langle\phi_j,f\rangle|^2.
\]
If \(f\) is smooth, then \(f\in H^s\) for every \(s\), and the spectral tail
converges to zero in every Sobolev norm. Sobolev embedding therefore gives
\(\Pi_\Lambda f\to f\) in \(C^k(X)\) for every \(k\).

The Schwartz kernel theorem identifies continuous operators
\(C^\infty(X)\to\Dcal'(X)\) with distributions on \(X\times X\). The preceding
convergence, together with the uniform Sobolev tail estimates below, implies
convergence of the corresponding kernels in the distribution topology. The
kernel of the limiting identity operator is
\(\delta_{\mathrm{diag}}\).
\end{proof}

\begin{theorem}[Quantitative spectral truncation]
\label{thm:spectral-error}
For \(r\geq0\), \(s\in\mathbb R\), and
\(f\in H^{s+r}(X)\),
\[
\|(I-\Pi_\Lambda)f\|_{H^s}
\leq
(1+\Lambda)^{-r/2}\|f\|_{H^{s+r}}.
\]
\end{theorem}

\begin{proof}
Writing \(f_j=\langle\phi_j,f\rangle\),
\begin{align*}
\|(I-\Pi_\Lambda)f\|_{H^s}^2
&=
\sum_{\lambda_j>\Lambda}
(1+\lambda_j)^s|f_j|^2\\
&\leq
(1+\Lambda)^{-r}
\sum_{\lambda_j>\Lambda}
(1+\lambda_j)^{s+r}|f_j|^2.
\end{align*}
Taking square roots gives the claim.
\end{proof}

The theorem says exactly what is finite and what is limiting. At finite
\(\Lambda\), \(K_\Lambda\) is the identity on the selected spectral subspace
and suppresses its orthogonal complement. The full delta appears only as the
subspaces exhaust \(L^2(X)\). Nothing in this theorem selects a physical value
of \(\Lambda\).

\section{Heat kernels are smoothing, not projections}

The heat semigroup gives another approximation to the identity:
\[
e^{-\tau\Delta}f
=\sum_j e^{-\tau\lambda_j}
\langle\phi_j,f\rangle\phi_j,
\qquad \tau>0.
\]
Its kernel is
\[
H_\tau(x,y)
=\sum_j e^{-\tau\lambda_j}
\phi_j(x)\overline{\phi_j(y)}.
\]

\begin{theorem}[Heat-kernel approximate identity]
\label{thm:heat}
Under Assumption~\ref{ass:elliptic},
\[
H_\tau\longrightarrow\delta_{\mathrm{diag}}
\quad\text{in }\Dcal'(X\times X)
\]
as \(\tau\downarrow0\). Moreover, for \(0\leq r\leq2\),
\(s\in\mathbb R\), and \(f\in H^{s+r}(X)\),
\[
\|(e^{-\tau\Delta}-I)f\|_{H^s}
\leq
\tau^{r/2}\|f\|_{H^{s+r}}.
\]
\end{theorem}

\begin{proof}
For \(u\geq0\) and \(0\leq\alpha\leq1\),
\[
0\leq1-e^{-u}\leq u^\alpha.
\]
Taking \(\alpha=r/2\) gives
\begin{align*}
\|(e^{-\tau\Delta}-I)f\|_{H^s}^2
&=
\sum_j(1+\lambda_j)^s
|1-e^{-\tau\lambda_j}|^2|f_j|^2\\
&\leq
\tau^r
\sum_j(1+\lambda_j)^s\lambda_j^r|f_j|^2\\
&\leq
\tau^r\|f\|_{H^{s+r}}^2.
\end{align*}
The operator convergence on smooth functions and the Schwartz kernel theorem
then give the distributional kernel limit.
\end{proof}

For small time and away from global complications, the familiar local
asymptotic form is
\[
H_\tau(x,y)
\sim
(4\pi\tau)^{-d/2}
\exp\!\left[-\frac{\dist(x,y)^2}{4\tau}\right]
\sum_{k\geq0}a_k(x,y)\tau^k.
\]
This explains the finite-width appearance of the heat kernel. It must not be
confused with spectral projection: \(e^{-\tau\Delta}\) has weights strictly
between zero and one on positive eigenspaces and is generally not idempotent.
The two constructions approach the same identity distribution through
different operator families.

\section{Normalized tubes around a constraint surface}

A hard constraint
\[
\delta(C(x))
\]
is not generally an identity kernel. It localizes an integral to a level set
and carries a geometric Jacobian.

Let
\[
g_\epsilon(z)
=(2\pi\epsilon^2)^{-m/2}
\exp\!\left(-\frac{|z|^2}{2\epsilon^2}\right),
\qquad z\in\mathbb R^m.
\]
The normalization is indispensable:
\(\int_{\mathbb R^m}g_\epsilon(z)\,\dd z=1\).

\begin{theorem}[Gaussian constraint-tube limit]
\label{thm:constraint-tube}
Let \(C:\mathbb R^n\to\mathbb R^m\), \(m\leq n\), be smooth, and let
\(f\in C_c(\mathbb R^n)\). Assume \(DC(x)\) has rank \(m\) on a neighborhood
of \(\supp(f)\). Define the normal Jacobian
\[
\Jcal_C(x)
=\sqrt{\det\!\bigl(DC(x)DC(x)^{\mathsf T}\bigr)}.
\]
Then
\[
\lim_{\epsilon\downarrow0}
\int_{\mathbb R^n}f(x)g_\epsilon(C(x))\,\dd x
=
\int_{C^{-1}(0)}
\frac{f(x)}{\Jcal_C(x)}
\,\dd\mathcal H^{n-m}(x).
\]
\end{theorem}

\begin{proof}
The coarea formula gives
\[
\int_{\mathbb R^n}f(x)g_\epsilon(C(x))\,\dd x
=
\int_{\mathbb R^m}g_\epsilon(z)F(z)\,\dd z,
\]
where
\[
F(z)
=
\int_{C^{-1}(z)}
\frac{f(x)}{\Jcal_C(x)}
\,\dd\mathcal H^{n-m}(x).
\]
Because \(C\) is a submersion near the compact support of \(f\), local
submersion coordinates and a finite partition of unity show that \(F\) is
continuous near \(z=0\). It is compactly supported. The family
\(\{g_\epsilon\}\) is an approximate identity on \(\mathbb R^m\), so the last
integral converges to \(F(0)\).
\end{proof}

This theorem gives a rigorous version of a finite admissibility tube. It also
shows why writing only
\(\exp[-|C|^2/(2\epsilon^2)]\) is incomplete: without the Gaussian
normalization, the integral generally tends to zero, and without the normal
Jacobian the limiting surface measure is wrong.

The result is local in constraint geometry. If zero is not a regular value,
the level set is singular and a different analysis is needed. If the
configuration space is infinite-dimensional, neither Lebesgue measure nor the
functional determinant follows from this finite-dimensional theorem.

\section{An exact projected commutator}

The identity kernel enters canonical field commutators. There is one precise
setting in which replacing the full identity by a projected kernel is exact.

\begin{theorem}[Compressed canonical commutation relation]
\label{thm:compressed-ccr}
Let \(\Hh\) be a complex one-particle Hilbert space, let
\(\Fock_s(\Hh)\) be its bosonic Fock space, and let \(P\) be an orthogonal
projection on \(\Hh\). On the usual finite-particle domain,
\[
[a(Pf),a^\dagger(Pg)]
=\langle Pf,Pg\rangle I
=\langle f,Pg\rangle I.
\]
If \(P\) is represented by a sufficiently regular integral kernel \(K_P\),
then
\[
\langle f,Pg\rangle
=
\int_{X\times X}
\overline{f(x)}K_P(x,y)g(y)
\,\dd V(x)\dd V(y).
\]
\end{theorem}

\begin{proof}
The canonical commutation relation on \(\Hh\) is
\[
[a(u),a^\dagger(v)]=\langle u,v\rangle I.
\]
Set \(u=Pf\) and \(v=Pg\). Since \(P=P^\ast=P^2\),
\[
\langle Pf,Pg\rangle=\langle f,Pg\rangle.
\]
The kernel expression is the definition of \(P\) as an integral operator.
\end{proof}

This is not a license to replace every
\(\delta(x-y)\) in a field theory by an arbitrary kernel. It says that if the
theory is explicitly compressed to \(\Ran(P)\), then the compressed creation
and annihilation operators satisfy the projected relation. Dynamics,
locality, covariance, positivity, and gauge symmetry must still be checked in
the compressed model.

\section{Fixed points: concentration and finite spread}

Fixed-point dynamics does not by itself force finite width. In fact, a
deterministic attracting fixed point naturally produces a delta limit.

\begin{proposition}[Pushforward toward a deterministic attractor]
\label{prop:fixed-point-delta}
Let \(S\) be a metric space, \(T:S\to S\) a measurable map, and \(x_\ast\in S\).
Suppose \(T^n x\to x_\ast\) for every \(x\) in a measurable basin \(B\). If
\(\mu\) is a probability measure supported in \(B\), then
\[
(T^n)_\#\mu\Longrightarrow\delta_{x_\ast}
\]
weakly.
\end{proposition}

\begin{proof}
For every bounded continuous \(h:S\to\mathbb R\),
\[
\int_S h\,\dd((T^n)_\#\mu)
=
\int_B h(T^n x)\,\dd\mu(x)
\longrightarrow
h(x_\ast)
\]
by bounded convergence.
\end{proof}

The width of the basin is irrelevant to this conclusion. If there are several
attracting basins, an initial ensemble can instead converge to a mixture of
point masses, with weights inherited from the initial measure. Neither case
selects a unique realized history from the limiting ensemble.

A nonzero stationary spread requires another ingredient. For example, the
declared stochastic equation
\[
\dd X_t=-\gamma X_t\,\dd t+\sqrt{2D}\,\dd W_t
\]
has stationary variance \(D/\gamma\) when \(\gamma,D>0\). That variance comes
from the balance between damping and continuing disturbance. It does not
follow from the fixed point or spectral gap alone, and it is not automatically
a detector width, coherence length, or MTT-selected constant. Finite-memory
driving changes the response again and must be analyzed in the declared
colored-noise model.

\section{Gauge fixing: what the Jacobian argument proves}

The Faddeev--Popov construction is often written schematically as
\[
1=
\int_{\mathcal G}
\delta\!\bigl(F(A^g)\bigr)
\det M_F(A^g)\,\dd g.
\]
Its finite-dimensional local content is a change-of-variables statement. Let
a Lie group \(G\) of dimension \(m\) act smoothly on a configuration manifold
\(\mathcal A\), and let \(F:\mathcal A\to\mathbb R^m\) be a gauge condition.
For fixed \(A\), consider the orbit map
\[
\Phi_A(g)=F(g\cdot A).
\]
If \(D_e\Phi_A\) is invertible and the chosen neighborhood contains exactly
one root of \(\Phi_A\), then the inverse function theorem and the ordinary
delta change-of-variables formula give a local identity of the form
\[
\int_G
\delta\!\bigl(\Phi_A(g)\bigr)
\left|\det D_g\Phi_A\right|
\,\dd g
=1.
\]
The determinant is the orbit-to-slice Jacobian. The normalized Gaussian tube
from Theorem~\ref{thm:constraint-tube} supplies a finite local approximation
to the slice constraint.

This local model explains the geometry but does not settle global gauge
fixing. In non-Abelian theories, one gauge condition can meet an orbit more
than once; these Gribov copies invalidate a naive global one-root identity
\cite{Gribov1978}. Infinite-dimensional path-integral measures and
determinants also require regularization. Ghost fields are a representation
of the determinant in the perturbative functional formalism, not proof that a
selected MTT kernel has supplied a global quotient.

\section{Measurement: a physical process with distinct layers}

For position on \(L^2(\mathbb R^d)\), the sharp observable is the
projection-valued measure
\[
(Q(B)\psi)(x)=\mathbf 1_B(x)\psi(x).
\]
The notation
\[
Q(\dd x)=|x\rangle\langle x|\,\dd x
\]
uses a distributional density. The generalized \(|x\rangle\) is not a vector
in \(L^2\), but every \(Q(B)\) is a bounded projection.

A normalized detector response \(g_\epsilon\) produces the smeared POVM
\[
E_\epsilon(B)
=
\int_{\mathbb R^d}
\left(\int_B g_\epsilon(y-x)\,\dd y\right)Q(\dd x).
\]
For a state \(\psi\), the outcome density is
\[
p_\epsilon=g_\epsilon*|\psi|^2.
\]
Normalized approximate identities give
\[
\|p_\epsilon-|\psi|^2\|_{L^1}\longrightarrow0,
\]
and the companion finite-resolution measurement paper proves a quantitative
\(W^{1,1}\) bound and constructs one compatible square-root instrument.

This example is a genuine finite-kernel bridge, but measurement contains more
than the POVM. The physical coupling, conditional state update, completion
into one outcome-bearing record, and later stabilization are distinct
layers. A detector-smearing theorem does not select one history. Conversely,
a discrete finite-dimensional PVM can be an ordinary bounded operator family
with no Dirac distribution at all \cite{DaviesLewis1970,Holevo2001}.

\section{Point sources, contacts, and conservation laws}

\subsection{Point sources}

An equation such as
\[
Lu=\delta_{x_0}
\]
defines a Green function or an ideal point-source response. Replacing the
source by a normalized profile \(\rho_\epsilon\) may be physically appropriate
when the source has finite extent:
\[
Lu_\epsilon=\rho_\epsilon,
\qquad
\rho_\epsilon\longrightarrow\delta_{x_0}.
\]
That is a source regularization, not necessarily a projection. The theory
must specify whether the point source is an exact mathematical probe, an
effective limit, or a claim about a physical emitter.

\subsection{Contact interactions}

A contact term such as \(\delta(x-y)\) idealizes a zero-range interaction.
Finite-range potentials can converge to contact models in declared scaling
limits, but the coupling may need dimension-dependent renormalization and the
operator domain can change. Merely substituting a smooth kernel does not prove
equivalence to the original contact theory.

\subsection{Momentum conservation}

Translation invariance gives the exact tempered-distribution identity
\[
\int_{\mathbb R^d}e^{iq\cdot x}\,\dd x
=(2\pi)^d\delta(q).
\]
Here the delta expresses exact Fourier orthogonality and conservation in an
infinite translation-invariant model. It is not evidence by itself for a
finite survivor basin.

If the interaction is restricted to a finite spacetime window \(W\), the
factor becomes
\[
\widehat{\mathbf 1_W}(q)
=
\int_W e^{iq\cdot x}\,\dd x.
\]
For rectangular windows this is a product of sinc profiles, and expanding
windows recover the conservation delta distributionally. The finite profile
therefore follows from a changed finite-window model. Whether the physical
system selects such a window is a separate question.

\section{Finite kernels do not automatically renormalize a theory}

A smoothing or finite-rank kernel can suppress high-frequency modes. It may
therefore be useful as a regulator or as the exact observable algebra of a
finite projected model. That observation is weaker than ultraviolet
completion.

An admissible replacement must address at least:
\begin{enumerate}
\item whether the kernel is selected or fitted;
\item whether gauge and Ward identities are preserved;
\item whether Lorentz covariance, causality, and unitarity survive;
\item whether the finite theory matches observed low-energy amplitudes;
\item whether the cutoff can be removed, or instead belongs to the exact
physical model; and
\item whether regulator dependence is controlled.
\end{enumerate}

Renormalization is not generally a repair of ``over-sharp projection.'' It is
a structured relation between bare descriptions, observables, scales, and
counterterms. A finite spectral algebra can make traces exact at its declared
cutoff, but connecting that algebra to a continuum QFT remains an additional
theorem.

\section{The MTT finite-kernel diagnostic}

The results above support a disciplined diagnostic rather than a universal
replacement rule.

\begin{definition}[Finite-kernel diagnostic]
For a downstream occurrence of a Dirac delta, the MTT finite-kernel diagnostic
asks:
\begin{enumerate}
\item Which mathematical role does the delta play?
\item Is there a finite projector, semigroup kernel, source profile, detector
response, constraint tube, or spacetime window appropriate to that role?
\item Which selected upstream object emits that finite representative?
\item In which topology does the representative converge to the delta?
\item What quantitative error certificate is available?
\item Which symmetries and operator identities survive at finite resolution?
\end{enumerate}
\end{definition}

This formulation is compatible with the projection-first program. The common
circle, lens filtration, nil/shear data, fixed-point maps, and finite q79
operators may help select particular finite objects. But those geometric
ingredients do not become a source theorem merely because a Gaussian or
spectral cutoff can be written down.

\subsection{What the present mathematics already supplies}

The current paper establishes four exact bridges:
\begin{enumerate}
\item finite spectral projectors converge to the diagonal identity
distribution with a Sobolev error;
\item heat kernels converge through a distinct smoothing family with a
Sobolev error;
\item normalized Gaussian constraint tubes converge to the coarea-weighted
surface distribution; and
\item compression to a declared one-particle subspace gives an exact
projected canonical commutator.
\end{enumerate}

These are reusable interfaces. A later MTT source theorem can point to one of
them and provide the missing finite object and scale without reproving the
analytic bridge.

\subsection{What remains open}

The selected MTT carrier has not yet been shown here to emit:
\begin{enumerate}
\item a universal spectral threshold \(\Lambda\) or heat time \(\tau\);
\item a detector response for arbitrary apparatus contexts;
\item a global non-Abelian gauge slice free of copy ambiguities;
\item a symmetry-preserving finite kernel for every QFT sector; or
\item a one-history completion law.
\end{enumerate}

The circle--lens--nil interpretation can organize candidate sources, but the
assignment of individual delta roles to those layers remains a research
proposal. In particular, there is no theorem here identifying a compact phase
circle with physical time or deriving a physical resolution scale from a
spectral gap.

\section{Status ledger}

\begin{center}
\begin{tabular}{p{0.20\linewidth}p{0.69\linewidth}}
\toprule
Status & Result\\
\midrule
Exact &
The diagonal delta is the Schwartz kernel of the identity\\
Exact &
Spectral projector kernels converge distributionally to the diagonal delta\\
Exact &
Spectral truncation has the stated \(H^{s+r}\to H^s\) error\\
Exact &
Heat kernels form a nonprojective approximate identity with the stated error\\
Exact &
Normalized Gaussian constraint tubes converge by the coarea formula\\
Exact &
Orthogonal compression gives the projected CCR kernel\\
Exact &
A deterministic attracting fixed point can concentrate an ensemble to a
delta\\
Standard input &
Detector POVMs, local Faddeev--Popov formalism, and distributional Fourier
conservation\\
Conditional &
Replacing a physical sharp object by a finite source, detector, window, or
kernel\\
Open &
Selected MTT scales, global gauge fixing, general instrument source,
one-history completion, and ultraviolet completion\\
\bottomrule
\end{tabular}
\end{center}

\section{Conclusion}

Dirac deltas are not all shadows of one hidden process. Their shared
distributional notation conceals several operator, source, constraint,
symmetry, and measurement roles. Once those roles are separated, a useful
projection-first insight survives in a rigorous form.

Spectral projector kernels and heat kernels both approach the identity
distribution, but one family is projective and the other is smoothing.
Normalized Gaussian tubes approach regular constraint surfaces with the
coarea Jacobian. Compressed field operators obey an exact projected
commutator. Finite detector responses and finite spacetime windows provide
further role-specific approximations.

For MTT, the productive question is therefore not ``which projection is every
delta hiding?'' It is ``which finite object does the selected geometry emit
for this role, and what proves the sharp limit without losing the required
physics?'' The analytic interfaces are now explicit. Selecting their physical
inputs remains the frontier.

\begin{thebibliography}{99}

\bibitem{BerlineGetzlerVergne1992}
N.~Berline, E.~Getzler, and M.~Vergne,
\emph{Heat Kernels and Dirac Operators},
Springer, Berlin, 1992,
doi:10.1007/978-3-642-58088-8.

\bibitem{DaviesLewis1970}
E.~B. Davies and J.~T. Lewis,
\emph{An Operational Approach to Quantum Probability},
Communications in Mathematical Physics \textbf{17} (1970) 239--260,
doi:10.1007/BF01647093.

\bibitem{EvansGariepy2015}
L.~C. Evans and R.~F. Gariepy,
\emph{Measure Theory and Fine Properties of Functions},
revised edition, CRC Press, Boca Raton, 2015,
doi:10.1201/b18333.

\bibitem{FaddeevPopov1967}
L.~D. Faddeev and V.~N. Popov,
\emph{Feynman Diagrams for the Yang--Mills Field},
Physics Letters B \textbf{25} (1967) 29--30,
doi:10.1016/0370-2693(67)90067-6.

\bibitem{Gribov1978}
V.~N. Gribov,
\emph{Quantization of Non-Abelian Gauge Theories},
Nuclear Physics B \textbf{139} (1978) 1--19,
doi:10.1016/0550-3213(78)90175-X.

\bibitem{Holevo2001}
A.~S. Holevo,
\emph{Statistical Structure of Quantum Theory},
Lecture Notes in Physics Monographs, vol.~67,
Springer, Berlin, 2001,
doi:10.1007/3-540-44998-1.

\end{thebibliography}

\end{document}
